\SourceLeaf{149}%
hängigen Erfahrungen, und kann über die eigentliche Normalzahl noch keinen viel sicherern Aufschluss geben, als die Erfahrung von einem einzelnen Jahre. Um das vergrösserte einem Durchschnittswerthe beizulegende Gewicht schätzen zu können, kommt es wesentlich darauf an, von wie vielen von einander ganz unabhängigen Gruppen die Erfahrungen hergenommen sind. Da nun die durchschnittliche Dauer eines Witwenthums gegen 20 Jahre beträgt, so würde bei empirischer Bestimmung der Normalzahl selbst ein vierzigjähriger Durchschnitt noch gar keine sehr sichere Bürgschaft für Elimination der Schwankungen gewähren. Dass vierzig Jahre nach einander die Witwenzahl beständig unter dem den allgemeinen Verhältnissen der Gesellschaft entsprechenden Normalwerthe bleibe, ist im Grunde eben so wenig für eine ganz ausserordentliche Erscheinung anzusehen, als wenn ein Pharaospieler zweimal nach einander gewinnt, und die Lehre, welche man hieraus ziehen muss, ist, dass von der andern Seite in jener Beziehung auch 40 ununterbrochen magere Jahre eben so leicht möglich sind, wie 40 ununterbrochen fette.

Das Zahlenverhältniss zwischen stehenden Ehen und Witwen, welches für die Professoren-Witwenkasse in \textsc{Kritter}’s Gutachten wie \(2:1\) vorausgesetzt ist, wird für Gesellschaften, welche verschiedenen Lebenskreisen angehören, ein sehr verschiedenes sein können. Vergleichen wir z. B. eine Witwenkassengesellschaft wie die unsrige mit der Gesammtheit aller stehenden Ehen und Witwen in einem ganzen Lande zunächst nur in Beziehung auf den allgemeinen [S. 140 oben] angegebenen Bestimmungsgrund, so sieht man leicht, dass in jener verhältnissmässig mehr Witwen gegen eine bestimmte Zahl von Ehen gerechnet werden müssen als in dieser. Bei der grossen Masse der Landeseinwohner fallen durchschnittlich die Verheirathungen in ein früheres Alter, und der Unterschied des Alters von Mann und Frau ist nach dem Durchschnittswerthe geringer, als bei Universitätsprofessoren\footnote{Aus unsrer Witwenkasse liegen mir die Data für das Alter nur von 62 Ehepaaren vor, wonach der durchschnittliche Unterschied 9 Jahre beträgt. Bei der Gesammtheit der Einwohner eines ganzen Landes wird man schwerlich auch nur einen halbsogrossen durchschnittlichen Unterschied annehmen können.}. Dazu kommt, dass in unsere Witwenkasse manche Mitglieder schon verheirathet eintreten, während in die Listen von einem ganzen Staate die sämmtlichen Ehepaare gleich von ihrer Verheirathung an eingerechnet werden. Allein die Wirkung dieser beiden Ursachen ist nur eine geringe im Vergleich zu den Einfluss eines andern Umstandes, welcher oben S. [140 ff.] bei der abstracten Behandlung der Sache noch bei Seite gesetzt wurde, nemlich, dass Abgang der Witwen nicht allein durch den Tod, sondern auch durch eine Wiederverheirathung erfolgen kann. In einer Witwenkasse, wo eine sich wieder verheirathende Witwe allen weitern Anspruch auf die Pension verliert, pflegen Wiederverheirathungen der Witwen selten vorzukommen, und namentlich zählt unsre Witwenkasse in dem ganzen Zeitraume ihres Bestehens nur einen einzigen Fall der Art; für ein ganzes Land hingegen nimmt man der Erfahrung zufolge an, dass aus dem ganzen Bestand der Witwen etwa der dreissigste Theil alljährlich durch Wiederverheirathung ausscheidet, und hiedurch wird das Verhältniss der stehenden Ehen zu den Witwen wesentlich abgeändert. Endlich hat auch in einem Lande, dessen Bevölkerung schon seit längerer Zeit im Zunehmen begriffen gewesen ist, diese Zunahme einen wesentlichen Einfluss auf das Verhältniss der \textit{coexistirenden} Ehen und Witwen, indem sich dann mehr stehende Ehen gegen Eine Witwe vorhanden werden, als ohne jenen Umstand da sein würden.

Als Erfahrungssatz wird gewöhnlich aufgestellt, dass für ein ganzes Land vier stehende Ehen gegen eine Witwe zu rechnen seien. Die von \textsc{Quetelet} für Belgien angegebenen Zahlen stimmen mit diesem Verhältniss fast genau überein. Für das ganze Königreich Hannover ergeben die Zählungen von 1833—1842 eine etwas kleinere Zahl, nemlich 3,74 stehende Ehen gegen eine Witwe; unterscheidet man aber die einzelnen Landestheile, so zeigen sich sehr grosse Ungleichheiten, es sind z. B.
